
\section{OA-08: Phase estimation with Fisher information}
\source{Week 6 Lecture 3, Quantum Optical Metrology; Mach--Zehnder slides.}
\subsection*{Step 1: identify the phase generator}
If a pure probe state is encoded by
\(\ket{\psi_\phi}=e^{-i\phi G}\ket{\psi_0}\),
the quantum Fisher information is
\[
\boxed{F_Q=4[\langle G^2\rangle-\langle G\rangle^2]}.
\]
For \(M\) independent identical probes and an unbiased
estimator the quantum Cram\'er--Rao inequality gives
\(\operatorname{Var}(\tilde\phi)\ge1/(MF_Q)\).

\subsection*{Step 2: calculate for a coherent optical phase}
The number operator generates an optical phase:
\(G=\hat n\). For \(\ket\alpha\),
\(\langle n\rangle=|\alpha|^2=\bar n\) and
\(\operatorname{Var}n=\bar n\), giving
\[
\boxed{F_Q=4\bar n,\quad
\Delta\tilde\phi\ge\frac1{2\sqrt{M\bar n}}}.
\]
This displays standard shot-noise scaling in total photon
resource \(M\bar n\). The factor of 2 depends on the
parameter-generator definition, not on a fundamental
change to the scaling.

\subsection*{Step 3: compare with actual measurement statistics}
For one photon in a balanced Mach--Zehnder interferometer,
write \(P_1=\cos^2(\phi/2)\),
\(P_2=\sin^2(\phi/2)\). The classical Fisher information is
\[
F=\sum_{j=1}^2\frac{(\partial_\phi P_j)^2}{P_j}.
\]
Since \(\partial_\phi P_1=-\tfrac12\sin\phi\),
\(\partial_\phi P_2=+\tfrac12\sin\phi\),
\[
F=\frac{\sin^2\phi}{4}
\left(\frac1{\cos^2(\phi/2)}+
\frac1{\sin^2(\phi/2)}\right)=\boxed{1}
\]
away from singular endpoints (with the continuous limit 1).
The one-photon phase bound for \(M\) trials is
\(\Delta\phi\ge1/\sqrt M\).
\checkmethod{The QFI is the optimal information available
from a \emph{state}; the classical Fisher information depends
on which measurement is actually made. State the resource count.}
